Sind keine textuellen Änderungen mehr zu erwarten, kann es an die Kontrolle
des Erscheinungsbildes gehen. Wer bis jetzt sauber gearbeitet hat, wird an
dieser Stelle nicht viel Arbeit haben. "`Dirty Hacks"' können einem nun jedoch
auf die Füße fallen. \TeX\ platziert in vielen Fällen einen Hinweis in der
Protokolldatei, wenn etwas nicht stimmt. Die meisten \LaTeX-Editoren
extrahieren die Warnungen im Protokollfile automatisch, stellen sie nach
Schweregrad farblich markiert dar und erlauben es, direkt an die entsprechende
Stelle im \TeX-Code zu springen. Soweit zutreffend, werden beispielhafte
Warnungen in den folgenden Unterabschnitten angegeben.
